\chapter{서\quad 론}
\label{ch:introduction}

\section{연구 배경 및 필요성}

본 템플릿은 가천대학교 일반대학원 학위논문 작성지침에 맞추어 구성된
\LaTeX\ 문서클래스 예시입니다. 본문은 11\,pt, 휴먼명조 계열, 줄간격
180\,\%\,$\sim$\,200\,\%로 설정되어 있으며, 문단 시작마다 2칸
들여쓰기가 적용됩니다.\footnote{각주는 9\,pt로 조판됩니다.}

두 번째 문단입니다. \texttt{indentfirst} 패키지를 사용하여 각 절의 첫 문단도
들여쓰기가 적용됩니다. 본문 내에서 그림과 표는 \verb|\ref{}| 로 상호참조할 수
있으며, 참고문헌은 \verb|\cite{}| 로 인용합니다~\cite{sample2026}.

\section{연구 목적}

본 연구의 목적은 다음과 같습니다.

\subsection{세부 목표 1}

여기에 첫 번째 세부 목표를 기술합니다.

\subsubsection{접근 방식}

(1) 번호 형식의 세부항 예시입니다.

\paragraph{단계 1} ① 번호 형식의 최하위 항목 예시입니다.

\subsection{세부 목표 2}

여기에 두 번째 세부 목표를 기술합니다.

\section{논문의 구성}

본 논문은 총 5개 장으로 구성됩니다. 제1장에서는 연구의 배경과 목적을
제시하고, 제2장에서는 관련 연구를 검토하며, 제3장에서는 제안 방법을,
제4장에서는 실험 결과를, 제5장에서는 결론을 제시합니다.
